
    %%%%%%%%%%%%%%%%%%%%%%%%%%%%%%%%%%%%%%%%%%%%%%%%%%%%%%%%%%%%%%%%%%%%%%%%%%%%%%%%
    %%% Classe do documento
    \documentclass[12pt, addpoints]{exam}
    \UseRawInputEncoding

    %%%%%%%%%%%%%%%%%%%%%%%%%%%%%%%%%%%%%%%%%%%%%%%%%%%%%%%%%%%%%%%%%%%%%%%%%%%%%%%%
    % preambulo.tex
    %
    % Arquivo de configuração de packages para uso o LaTeX, conforme minhas
    % preferências e modelos pessoais.
    %
    % Para maiores informações, visite:
    %    
    %
    % NÃO ALTERE SE NÃO SOUBER O QUE ESTÁ FAZENDO!
    %%%%%%%%%%%%%%%%%%%%%%%%%%%%%%%%%%%%%%%%%%%%%%%%%%%%%%%%%%%%%%%%%%%%%%%%%%%%%%%%


    %%%%%%%%%%%%%%%%%%%%%%%%%%%%%%%%%%%%%%%%%%%%%%%%%%%%%%%%%%%%%%%%%%%%%%%%%%%%%%%%
    %%% Estruturas de controle:
    \input{static/utils/controles}

    %%%%%%%%%%%%%%%%%%%%%%%%%%%%%%%%%%%%%%%%%%%%%%%%%%%%%%%%%%%%%%%%%%%%%%%%%%%%%%%%
    %%% Compilação condicional em PDF:
    \input{static/utils/pdf}

    %%%%%%%%%%%%%%%%%%%%%%%%%%%%%%%%%%%%%%%%%%%%%%%%%%%%%%%%%%%%%%%%%%%%%%%%%%%%%%%%
    %%% Configurações de layout da página:
    \input{static/utils/layout}

    %%%%%%%%%%%%%%%%%%%%%%%%%%%%%%%%%%%%%%%%%%%%%%%%%%%%%%%%%%%%%%%%%%%%%%%%%%%%%%%%
    %%% Configurações para autores:
    \input{static/utils/autores}

    %%%%%%%%%%%%%%%%%%%%%%%%%%%%%%%%%%%%%%%%%%%%%%%%%%%%%%%%%%%%%%%%%%%%%%%%%%%%%%%%
    %%% Cabeçalho e rodapé:
    %%% Atenção! É MUITO DIFÍCIL ajustar apropriadamente os running headers
    %%% e running footers da primeira vez. Você pode confiar no LaTeX e deixar que
    %%% ele faça o ajuste automaticamente ou pode quebrar a cabeça e
    %%% tentar melhorar algo que o LaTeX já faz muito bem, mesmo que não fique
    %%% exatamente do jeito que você quer. O que você prefere? Se quiser ajustar
    %%% manualmente, descomente a linha abaixo e faça as configurações no arquivo
    %%% "utils/headings.tex".
    %\input{static/utils/headings}

    %%%%%%%%%%%%%%%%%%%%%%%%%%%%%%%%%%%%%%%%%%%%%%%%%%%%%%%%%%%%%%%%%%%%%%%%%%%%%%%%
    %%% Configuração para as fontes do documento:
    %%% Se você quiser usar fontes próprias ou do sistema, precisará ajustar as
    %%% configurações no arquivo "utils/fontes.tex".
    %%% ATENÇÃO: por padrão eu utilizo XeLaTeX com fontes proprietárias projetadas
    %%% por Matthew Butterick, adquiridas em https://mbtype.com, que não podem ser
    %%% distribuídas devido à licença de utilização. Se você usar XeLaTeX, você
    %%% DEVERÁ OBRIGATORIAMENTE ajustar as fontes no arquivo "utils/fontes.tex".
    \input{static/utils/fontes}

    %%%%%%%%%%%%%%%%%%%%%%%%%%%%%%%%%%%%%%%%%%%%%%%%%%%%%%%%%%%%%%%%%%%%%%%%%%%%%%%%
    %%% Sumário:
    \input{static/utils/toc}

    %%%%%%%%%%%%%%%%%%%%%%%%%%%%%%%%%%%%%%%%%%%%%%%%%%%%%%%%%%%%%%%%%%%%%%%%%%%%%%%%
    %%% Matemática:
    \input{static/utils/matematica}

    %%%%%%%%%%%%%%%%%%%%%%%%%%%%%%%%%%%%%%%%%%%%%%%%%%%%%%%%%%%%%%%%%%%%%%%%%%%%%%%%
    %%% Notas de rodapé:
    \input{static/utils/footnotes}

    %%%%%%%%%%%%%%%%%%%%%%%%%%%%%%%%%%%%%%%%%%%%%%%%%%%%%%%%%%%%%%%%%%%%%%%%%%%%%%%%
    %%% Referências cruzadas, links, citações:
    \input{static/utils/crossrefs}

    %%%%%%%%%%%%%%%%%%%%%%%%%%%%%%%%%%%%%%%%%%%%%%%%%%%%%%%%%%%%%%%%%%%%%%%%%%%%%%%%
    %%% Configurações para os hiperlinks em PDF:
    \input{static/utils/pdfconfig}

    %%%%%%%%%%%%%%%%%%%%%%%%%%%%%%%%%%%%%%%%%%%%%%%%%%%%%%%%%%%%%%%%%%%%%%%%%%%%%%%%
    %%% Glossário e índice remissivo:
    \input{static/utils/glossindex}

    %%%%%%%%%%%%%%%%%%%%%%%%%%%%%%%%%%%%%%%%%%%%%%%%%%%%%%%%%%%%%%%%%%%%%%%%%%%%%%%%
    %%% Referências bibliográficas:
    \input{static/utils/bibliografia}

    %%%%%%%%%%%%%%%%%%%%%%%%%%%%%%%%%%%%%%%%%%%%%%%%%%%%%%%%%%%%%%%%%%%%%%%%%%%%%%%%
    %%% Cores:
    \input{static/utils/cores}

    %%%%%%%%%%%%%%%%%%%%%%%%%%%%%%%%%%%%%%%%%%%%%%%%%%%%%%%%%%%%%%%%%%%%%%%%%%%%%%%%
    %%% Computação:
    \input{static/utils/computacao}

    %%%%%%%%%%%%%%%%%%%%%%%%%%%%%%%%%%%%%%%%%%%%%%%%%%%%%%%%%%%%%%%%%%%%%%%%%%%%%%%%
    %%% Gráficos:
    \input{static/utils/graficos}

    %%%%%%%%%%%%%%%%%%%%%%%%%%%%%%%%%%%%%%%%%%%%%%%%%%%%%%%%%%%%%%%%%%%%%%%%%%%%%%%%
    %%% Ícones:
    \input{static/utils/icones}

    %%%%%%%%%%%%%%%%%%%%%%%%%%%%%%%%%%%%%%%%%%%%%%%%%%%%%%%%%%%%%%%%%%%%%%%%%%%%%%%%
    %%% Blocos:
    \input{static/utils/blocos}

    %%%%%%%%%%%%%%%%%%%%%%%%%%%%%%%%%%%%%%%%%%%%%%%%%%%%%%%%%%%%%%%%%%%%%%%%%%%%%%%%
    %%% Boxes coloridos:
    \input{static/utils/colorbox}

    %%%%%%%%%%%%%%%%%%%%%%%%%%%%%%%%%%%%%%%%%%%%%%%%%%%%%%%%%%%%%%%%%%%%%%%%%%%%%%%%
    %%% Tabelas:
    \input{static/utils/tabelas}

    %%%%%%%%%%%%%%%%%%%%%%%%%%%%%%%%%%%%%%%%%%%%%%%%%%%%%%%%%%%%%%%%%%%%%%%%%%%%%%%%
    %%% Ambientes de listas:
    \input{static/utils/listas}

    %%%%%%%%%%%%%%%%%%%%%%%%%%%%%%%%%%%%%%%%%%%%%%%%%%%%%%%%%%%%%%%%%%%%%%%%%%%%%%%%
    %%% Ambientes floats e similares:
    \input{static/utils/floats}

    %%%%%%%%%%%%%%%%%%%%%%%%%%%%%%%%%%%%%%%%%%%%%%%%%%%%%%%%%%%%%%%%%%%%%%%%%%%%%%%%
    %%% Ajustes para a classe EXAM:
    \input{static/utils/exam}

    %%%%%%%%%%%%%%%%%%%%%%%%%%%%%%%%%%%%%%%%%%%%%%%%%%%%%%%%%%%%%%%%%%%%%%%%%%%%%%%%
    %%% Ajustes para textos:
    \input{static/utils/texto}

    %%%%%%%%%%%%%%%%%%%%%%%%%%%%%%%%%%%%%%%%%%%%%%%%%%%%%%%%%%%%%%%%%%%%%%%%%%%%%%%%
    %%% Ambientes, aliases, comandos e customizações em geral para serem
    %%% utilizadas nos documentos. Defina aqui qualquer customização específica
    %%% para seus documentos!
    \input{static/utils/customizacoes}

    %%%%%%%%%%%%%%%%%%%%%%%%%%%%%%%%%%%%%%%%%%%%%%%%%%%%%%%%%%%%%%%%%%%%%%%%%%%%%%%%
    %%% Ajuste do layout, espaçamento de linhas e etc.:
    \geometry{a4paper, portrait, left=2.0cm, right=2.0cm, top=2.0cm, bottom=2.0cm}
    \singlespacing

    %%%%%%%%%%%%%%%%%%%%%%%%%%%%%%%%%%%%%%%%%%%%%%%%%%%%%%%%%%%%%%%%%%%%%%%%%%%%%%%%
    %%% Configurações de cabeçalho e rodapé (não use fancyhdr pois ocorre conflito):
    \pagestyle{headandfoot}
    \runningheadrule
    \firstpageheader{}{}{}
    \runningheader{Sem disciplina}{}{Avaliação}
    \firstpagefooter{}{}{Página \thepage\ de \numpages}
    \runningfooter{}{}{Página \thepage\ de \numpages}
    \runningfootrule

    %%%%%%%%%%%%%%%%%%%%%%%%%%%%%%%%%%%%%%%%%%%%%%%%%%%%%%%%%%%%%%%%%%%%%%%%%%%%%%%%
    %%% Configurações para as propriedades do PDF:
    \hypersetup{
    hidelinks,           % Comente para web, descomente para impressão
    colorlinks=true,      % True para web, False para impressão
    pdftitle=Nova Prova: Avaliação,
    pdfauthor=Professor,
    pdfsubject=Sem disciplina,
    pdfkeywords=['Sem tag'],
    pdfinfo={
        CreationDate={}, % Ex.: D:AAAAMMDDHH24MISS
        ModDate={}       % Ex.: D:AAAAMMDDHH24MISS
    }
    }
    \input{static/utils/pdfconfig}

    %%%%%%%%%%%%%%%%%%%%%%%%%%%%%%%%%%%%%%%%%%%%%%%%%%%%%%%%%%%%%%%%%%%%%%%%%%%%%%%%
    %%% Começa o documento
    \begin{document}

    %%%%%%%%%%%%%%%%%%%%%%%%%%%%%%%%%%%%%%%%%%%%%%%%%%%%%%%%%%%%%%%%%%%%%%%%%%%%%%%%
    %%% Folha de instruções padronizadas da para a prova
    \pdfbookmark[1]{Instruções}{instrucoes}
    \begin{figure}[H]
        \begin{center}
            \includegraphics[scale=0.3]{{static/imgs/logo-.png}}
        \end{center}	
    \end{figure}

    \vspace{-1cm}

    \begin{table}[H]
        \centering
        \begin{tabular}{|llll|}
            \hline
            \multicolumn{2}{|l|}{\textbf{Disciplina:} Sem disciplina} & 
            \multicolumn{1}{l|}{\multirow{3}{*}{\textbf{\makecell{Nota:\\ \,\\ \,}}}} & 
            \multirow{3}{*}{\textbf{\makecell{Coordenador:\\ \,\\ \,}}} \\ 
            \cline{1-2}
            \multicolumn{2}{|l|}{\textbf{Professor:} Professor \hspace{4.72cm} } & 
            \multicolumn{1}{l|}{} & \\ 
            \cline{1-2}
            \multicolumn{2}{|l|}{\textbf{Aluno:}} & 
            \multicolumn{1}{l|}{} & \\ 
            \hline
            \multicolumn{1}{|l|}{\textbf{Turma:} \hspace{6.3cm} } & 
            \multicolumn{1}{l|}{\textbf{Semestre:}} & 
            \multicolumn{2}{l|}{\textbf{Valor:} 10 pontos} \\ 
            \hline
            \multicolumn{1}{|l|}{\textbf{Data:}} & 
            \multicolumn{3}{l|}{\textbf{Avaliação:}} \\ 
            \hline
        \end{tabular}
    \end{table}

    \begin{center}
        \fbox{\setlength\fboxsep{0.3cm} \fbox{\parbox{15.4cm}{
            \begin{center}
                \textbf{INSTRUÇÕES PARA A REALIZAÇÃO DESTA AVALIAÇÃO:}
            \end{center}
            \begin{itemize}[leftmargin=*]
                \item Esta é a primeira avaliação da disciplina Sem disciplina, 
                    e engloba o conteúdo referente Sem tag.
                \item Esta avaliação contém 15 questões objetivas, \textbf{todas obrigatórias}.
                \item \textbf{Leia atentamente cada questão!} As questões objetivas terão uma,
                    e apenas uma única, resposta a ser marcada.
                \item \textbf{Desligue o celular} antes de começar. A avaliação é
                    \textbf{individual} e \textbf{sem consulta}.
                \item As questões podem ser respondidas, na prova, com lápis ou caneta. Ao final
                    da prova \textbf{{VOCÊ DEVERÁ PREENCHER O GABARITO COM CANETA
                    DE COR PRETA}} \textbf{\underline{OU AZUL}} para que seja feita a correção
                    automatizada através da leitura do padrão de respostas marcado no
                    gabarito.
                \item \textbf{CUIDADO ao preencher o gabarito} pois eles são \textbf{INDIVIDUAIS
                    e NOMEADOS} (cada estudante tem um gabarito próprio específico, com um
                    código de identificação único). \textbf{Questões rasuradas são anuladas}
                    pelo software de leitura do gabarito.
                \item Ao final da prova, \textbf{DEVOLVA A PROVA E O GABARITO} para o professor.
                \item Siga todas as normas de \textbf{Integridade Acadêmica} da disciplina.
                    Alunos que forem flagrados com qualquer espécie de ``cola'' ou trocando
                    informações com outros alunos terão suas avaliações recolhidas, as notas
                    zeradas, e a situação será encaminhada para a coordenação para a aplicação
                    das penalidades previstas pela universidade.
                \item Com 90 minutos de avaliação você terá tempo de até 6,min por questão,
                    em média.
                \item Boa avaliação!
            \end{itemize}
            \vspace{2cm}
        }}}
    \end{center}
    \newpage

    %%%%%%%%%%%%%%%%%%%%%%%%%%%%%%%%%%%%%%%%%%%%%%%%%%%%%%%%%%%%%%%%%%%%%%%%%%%%%%%%
    %%% Inicia o bloco de questões:
    \begin{questions}

    \newpage
    \question

A derivada da função \( f(x) = x^x \) é igual a:
\begin{choices}
\choice \( x^x \ln(x) \)
\choice \( x^x \)
\choice \( x x^{x-1} \)
\choice \( x^x (\ln(x) + x) \)
\choice \( x^x (\ln(x) + 1) \)
\end{choices}

\question

Você estava ouvindo a discussão de dois colegas, o João e a Maria. Enquanto o

João afirmava que o conceito de banco de dados (BD) é extremamente diferente do

conceito de sistema de gerenciamento de bancos de dados (SGBD), a Maria

discordava e afirmava que o banco de dados nada mais é do que um componente do

SGBD. Nesse momento o professor entrou na sala, percebeu a discussão e falou

que:
\begin{choices}
\choice Não é possível saber quem estava certo ou errado pois não definiram qual o modelo de dados que estava sendo considerado na discusão.
\choice Os dois estavam errados pois não são conceitos separados nem componentes, são conceitos sinônimos.
\choice O João estava certo já que o banco de dados é, realmente, um conceito diferente do conceito do sistema de gerenciamento de bancos de dados, e existe de forma independente.
\choice Os dois estavam certos pois o conceito, na prática, depende da situação específica que se está considerando.
\choice A Maria estava certa já que o banco de dados é, realmente, um componente interno do sistema de gerenciamento de bancos de dados e não existe de forma independente.
\end{choices}

\question

(ENADE 2014) Considere o seguinte banco de dados ilustrado no diagrama relaconal

abaixo:

\begin{figure}[H]

  \centering

  \caption{Estados e fornecedores}

  \label{fig:estados}

  %\fbox{

    \includegraphics[width=0.5\linewidth]{static/imgs/defaul.png}

  %}\\

  %\footnotesize{Fonte: xxxx.}

\end{figure}

A expressão SQL que obtém os nomes dos estados para os quais não há fornecedores

cadastrados é:


\begin{choices}
\choice \begin{verbatim}SELECT e.nome

FROM estado e,

FROM fornecedor f

WHERE e.nome = f.uf;

\end{verbatim}


\choice \begin{verbatim}SELECT e.nome

FROM estado e

WHERE e.uf IN (SELECT f.uf

               FROM fornecedor f);

\end{verbatim}


\choice \begin{verbatim}SELECT e.nome

FROM estado e

WHERE e.uf NOT IN (SELECT f.uf

                   FROM fornecedor f);

\end{verbatim}


\choice \begin{verbatim}SELECT e.nome

FROM estado e,

FROM fornecedor f

WHERE e.uf = f.uf;

\end{verbatim}


\choice \begin{verbatim}SELECT e.uf

FROM estado e

WHERE e.nome NOT IN (SELECT f.uf

                     FROM fornecedor f);

\end{verbatim}
\end{choices}

\question

Na criptografia com chave pública:


\begin{choices}
\choice O sigilo é obtido através da codificação com a chave privada do remetente e decifragem com a chave pública do destinatário.
\choice O sigilo é obtido através da codificação com a chave privada do destinatário e decifragem com a chave pública do destinatário.
\choice Para assinar digitalmente uma mensagem codifica-se a mesma com a chave pública do destinatário e esta é decifrada com a chave privada do destinatário.
\choice Para assinar digitalmente uma mensagem codifica-se a mesma com a chave pública do remetente e esta é decifrada com a chave privada do destinatário.
\choice O sigilo é obtido através da codificação com a chave pública do destinatário e decifragem com a chave privada do destinatário.
\end{choices}

\question

(Adaptado do ENADE 2005) Avalie o seguite diagrama relacional de um banco de dados:

\begin{figure}[H]

  \centering

  \caption{Banco de dados de peças e fornecedores}

  \label{fig:pecas}

  %\fbox{

    \includegraphics[width=0.5\linewidth]{static/imgs/defaul.png}

  %}\\

  %\footnotesize{Fonte: xxxx.}

\end{figure}

Assinale a opção que apresenta um comando SQL que permite obter uma lista que

contenha o nome de cada fornecedor que tenha fornecido alguma peça, o código da

peça fornecida, a descrição dessa peça e a quantidade fornecida da referida

peça.
\begin{choices}
\choice \begin{verbatim}SELECT * FROM pecas INNER JOIN fornecedores INNER JOIN fornecimentos;\end{verbatim} 
\choice \begin{verbatim}SELECT f2.nome, p.codigo, p.descricao, f.quantidade FROM pecas INNER JOIN fornecimentos f ON (p.codigo = f.cod_peca) INNER JOIN fornecedores f2 ON (f.cod_forn = f2.cod_forn);\end{verbatim} 
\choice \begin{verbatim}SELECT * FROM pecas INNER JOIN fornecimentos f ON (p.codigo = f.cod_peca) INNER JOIN fornecedores f2 ON (f.cod_forn = f2.cod_forn) ORDER BY f2.nome;\end{verbatim} 
\choice \begin{verbatim}SELECT DISTINCT p.nome, p.codigo, p.descricao, f.quantidade FROM pecas p INNER JOIN fornecimentos f ON (p.codigo = f.cod_peca);\end{verbatim} 
\choice \begin{verbatim}SELECT nome, codigo, descricao, quantidade FROM pecas INNER JOIN FORNECEDORES INNER JOIN FORNECIMENTOS;\end{verbatim} 
\end{choices}

\question

Qual das seguintes igualdades é verdadeira?



\begin{enumerate}

    \item[(I)] \( n^2 = \mathcal{O}(n^3) \)

    \item[(II)] \( 2 \cdot n + 1 = \mathcal{O}(n^2) \)

    \item[(III)] \( n^3 = \mathcal{O}(n^2) \)

    \item[(IV)] \( 3 \cdot n + 5 \cdot n \log n = \mathcal{O}(n) \)

    \item[(V)] \( \log n + \sqrt{n} = \mathcal{O}(n) \)

\end{enumerate}



As opções para a resposta correta são:
\begin{choices}
\choice somente II, III e IV
\choice somente I, II e V
\choice somente I e II
\choice somente III, IV e V
\choice somente I, III e IV
\end{choices}

\question

Seja a árvore binária abaixo a representação de um espaço de estados para um problema \( p \), em que o estado inicial é \( a \), e \( i \) e \( f \) são estados finais.



\begin{figure}[H]

    \centering

    \includegraphics[width=0.5\linewidth]{static/imgs/defaul.png}

    \caption{Questão 59}

    \label{fig:enter-label}

\end{figure}



Um algoritmo de busca em largura-primeiro forneceria a seguinte sequência de estados como primeira alternativa a um caminho-solução para o problema \( p \):
\begin{choices}
\choice \( a \ b \ c \ d \ e \ f \)
\choice \( a \ b \ e \ i \)
\choice \( a \ b \ d \ h \ e \ i \)
\choice \( a \ c \ f \)
\choice \( a \ b \ d \ e \ f \)
\end{choices}

\question

Ao analisar um novo sistema de gerenciamento de bancos de dados que surgiu no

mercado, você percebe que a entrada às operações do modelo de dados são tabelas

mas as saídas nunca são tabelas, são sempre outras estruturas. O vendedor desse

SGBD afirma que ele é um SGBD relacional. Pode-se afirmar que:


\begin{choices}
\choice Esse SGBD não é relacional pois no modelo relacional a única estrutura percebida pelos usuários são tabelas.
\choice Esse SGBD é relacional pois consegue trabalhar com tabelas para a entrada das operações do modelo de dados. A saída não é importante.
\choice Nenhuma das alternativas anteriores.
\choice Esse SGBD não é relacional pois, para ser verdadeiramente relacional, as operações do modelo de dados não podem trabalhar recebendo tabelas como entrada.
\choice Esse SGBD é relacional pois o vendedor afirmou que é.
\end{choices}

\question

Dentre as definições a seguir, ligadas ao conceito de visões do modelo relacional, qual delas é \textbf{INCORRETA}?


\begin{choices}
\choice Uma visão relacional é uma relação virtual que nunca é materializada.
\choice Programas aplicativos do banco de dados podem ser executados sobre visões de relações da base de dados.
\choice Uma visão relacional é uma relação virtual, derivada de relações base a partir da especificação de operações da álgebra relacional.
\choice O gerenciamento de visões envolve a conversão da consulta do usuário sobre as visões para a consulta sobre as relações base.
\choice Uma visão é útil por representar uma percepção particular do banco de dados, compartilhado por muitos aplicativos.
\end{choices}

\question

(ENADE 2017) No modelo relacional de dados, uma junção entre tabelas cria uma

outra tabela derivada de duas ou mais tabelas de acordo com os critérios

especificados para a junção, de modo similar às regras da teoria dos

conjuntos. Nos conjuntos a seguir, as áreas hachuradas representam os resultados

das consultas SQL em que ``id' é uma chave primária e ``fk' é uma chave

estrangeira:

\begin{figure}[H]

  \centering

  \includegraphics[width=0.5\linewidth]{static/imgs/defaul.png}

\end{figure}

Assinale a opção em que a declaração SQL é equivalente ao conjunto apresentado a

seguir:

\begin{figure}[H]

  \centering

  \includegraphics[width=0.5\linewidth]{static/imgs/defaul.png}

\end{figure}


\begin{choices}
\choice \begin{verbatim}SELECT *

FROM tabela1 t1

WHERE NOT EXISTS (SELECT 1

                  FROM tabela1 t2

                  WHERE t2.id = tk.fk);

\end{verbatim}


\choice \begin{verbatim}SELECT *

FROM tabela1 t1 WHERE EXISTS (SELECT 1

                              FROM tabela2 t2

                              WHERE t2.id = t1.fk);

\end{verbatim}


\choice \begin{verbatim}SELECT *

FROM tabela1 t1

RIGHT JOIN tabela2 t2 ON (t1.id = t2.fk)

WHERE t2.fk <> NULL;

\end{verbatim}


\choice \begin{verbatim}SELECT *

FROM tabela1 t1

RIGHT JOIN tabela2 t2 ON (t1.id = t2.fk)

WHERE t1.id IS NULL;

\end{verbatim}


\choice \begin{verbatim}SELECT *

FROM tabela1 t1

INNER JOIN tabela2 t2 ON (t1.id = t2.fk);

\end{verbatim}


\end{choices}

\question

\textit{Starvation} ocorre quando:
\begin{choices}
\choice Dois ou mais processos são forçados a acessar dados críticos alternando estritamente entre eles.
\choice O processo tenta mas não consegue acessar uma variável compartilhada.
\choice A prioridade de um processo é ajustada de acordo com o tempo total de execução do mesmo.
\choice Pelo menos um evento espera por um evento que não vai ocorrer.
\choice Pelo menos um processo é continuamente postergado e não executa.
\end{choices}

\question

A interposição de um circuito de memória cache entre o processador e a memória principal (RAM)
\begin{choices}
\choice aumenta o tráfego de instruções e/ou dados entre memória e disco
\choice aumenta o tráfego de instruções e/ou dados no barramento de memória
\choice permite acessos concorrentes à memória RAM
\choice diminui o tráfego de instruções e/ou dados entre memória e disco
\choice diminui o tráfego de instruções e/ou dados no barramento de memória
\end{choices}

\question

Algoritmos distribuídos podem usar passagem de "token" por um anel lógico para implementar exclusão mútua ou ordenação global de mensagens. Nesses algoritmos, apenas o processo que possui o "token" tem a permissão de usar um recurso compartilhado ou numerar mensagens, por exemplo. Considerando o conceito acima, podemos afirmar que: 
\begin{choices}
\choice para usar essa abordagem, os computadores precisam estar conectados em uma rede com topologia em 
\choice nessa abordagem é impossível evitar a geração espontânea de vários "tokens" mesmo em sistemas livres de falhas 
\choice a abordagem é pouco robusta pois a perda do "token" por um processo provoca o bloqueio do algoritmo distribuído que a usa 
\choice a abordagem é adequada apenas para sistemas onde possa ser controlado o tempo que cada computador permanece com o 
\choice a abordagem deve tratar no mínimo dois tipos de defeitos: perda do "token" e colapso de processos 
\end{choices}

\question

Qual o valor do atributo \( E.val \) após a análise da expressão "4 / 2 / 2" para o esquema de tradução a seguir?



\begin{itemize}

    \item \( E \rightarrow T / E1 \ \{ E.val = T.val / E1.val \} \)

    \item \( E \rightarrow T \ \{ E.val = T.val \} \)

    \item \( T \rightarrow \text{digito} \ \{ T.val = \text{val(digito)} \} \)

\end{itemize}


\begin{choices}
\choice 1
\choice 8
\choice 4
\choice 3
\choice 2
\end{choices}

\question

Um banco de dados (BD) é uma coleção logicamente coerente de dados relacionados

e persistentes. Em relação às características de um banco de dados, assinale a

alternativa verdadeira:
\begin{choices}
\choice Não precisa de dados atuais e verdadeiros para ser útil (para representar corretamente a realidade atual).
\choice Geralmente é projetado e usado para finalidades gerais, sendo extremamente genéricos e adaptáveis para qualquer situação.
\choice Não precisam ser obrigatoriamente computadorizados, existem bancos de dados manuais e/ou em papel.
\choice Armazena dois grandes ``tipos`` de dados: os dados do usuário final e os metadados, sendo que os metadados ficam armazenados nas mesmas tabelas que os dados dos usuários e isso nos permite saber a estrutura dos bancos de dados.
\choice Representa um aspecto (ou parte) do mundo real, parte essa chamada de ``catálogo' ou ``dicionário de dados``.
\end{choices}



    %%%%%%%%%%%%%%%%%%%%%%%%%%%%%%%%%%%%%%%%%%%%%%%%%%%%%%%%%%%%%%%%%%%%%%%%%%%%%%%%
    %%% Finaliza o bloco de questões:
    \end{questions}
    \end{document}
    